%=== APPENDICES ===
\begin{appendices}
\label{cha:appendices}

\chapter{Implemented OCP Configuration}
\markboth{Appendix A}{}

\begin{spacing}{1.5}
\begin{table}[ht]
\centering
\caption{NMPC configuration used as the dissertation baseline.}
\begin{tabular}{@{}ll@{}}
\toprule
State dimension & 13 \\
Input dimension & 4 \\
Prediction stages & $N_c=20$ \\
Sampling interval & $\Delta t=\SI{0.05}{s}$ \\
Prediction horizon & \SI{1.0}{s} \\
Integrator & Explicit RK4 \\
NLP method & Full SQP with merit backtracking \\
QP solver & Partial-condensing HPIPM \\
Hessian & Gauss--Newton \\
\bottomrule
\end{tabular}
\end{table}

\begin{table}[ht]
\centering
\caption{MHE configuration used as the dissertation baseline.}
\small
\begin{tabularx}{\textwidth}{@{}p{0.31\textwidth} X@{}}
\toprule
Augmented state dimension & 14 \\
Process-noise dimension & 13 \\
Known-input dimension & 7 (thrust, torque and CoM geometry) \\
Estimation stages & $N_e=20$ \\
Sampling interval & $\Delta t=\SI{0.1}{s}$ \\
Moving horizon & \SI{2.0}{s} \\
Mass interval & \SIrange{0.2}{5.0}{kg} (numerical only) \\
Mass arrival weight & $Q_{0,m}=10^{-6}$ \\
Initial iterate & $\operatorname{clip}(T_{\mathrm{phys}}/g,0.2,5.0)$ \\
Pre-offboard gate & $z>\SI{0.5}{m}$; 20 frames, $\sigma_m<\SI{0.02}{kg}$ \\
Payload prediction & $\hat m_p=(\hat m_t-m_b)^+$ \\
Online inertia tracking & \texttt{DJ\_TRACK\_MEST=1}; ratchet selectable \\
Inertia-map safety bounds & $m_{p,\mathrm{floor}}=\SI{0.15}{kg}$; $m_{p,\max}=\SI{0.6}{kg}$ \\
Gain initialisation envelope & \SI{0.5}{kg} platform rating; not supplied to MHE \\
Event schedule & M0: pre-event $R/Q$ scale $10^{-4}$ \\
Step trigger & 0.3-s half-windows; \SI{1.0}{N}; 2-frame persistence \\
Drop release threshold & signed step below \SI{-2.0}{N} \\
Drop state action & release geometry; never reset mass \\
Baseline switches & no-mass-prior, pre-offboard, coupled MHE and online $\dJ$ on \\
NLP/QP method & SQP / partial-condensing HPIPM \\
\bottomrule
\end{tabularx}
\end{table}
\end{spacing}

\clearpage
\chapter{Experiment Record Template}
\markboth{Appendix B}{}

\begin{spacing}{1.5}
Each final run should store the following metadata with the log:
\begin{enumerate}
    \item Git revision, build timestamp and generated-solver checksum;
    \item vehicle mass/inertia calibration and fixed actuator limits;
    \item payload mass for offline scoring, attachment geometry and proof that payload truth is absent from the adaptive path;
    \item trajectory parameters, grasp/release phase and random seed;
    \item MHE/NMPC horizons, weights, bounds, no-prior switches and solver options;
    \item estimator/control topic rates and measured end-to-end latency;
    \item termination reason, solver status sequence and constraint violations.
\end{enumerate}

A final results table should contain one row per run rather than only an average. This permits the aggregated figures to be regenerated and prevents failed runs from disappearing during manual selection.
\end{spacing}

\chapter{Pre-Submission Completion Checklist}
\markboth{Appendix C}{}

\begin{spacing}{1.5}
\begin{enumerate}
    \item Replace candidate, programme, supervisor and acknowledgement placeholders.
    \item Confirm the declaration statements against the current NTU MAE wording.
    \item Set \texttt{\string\showdraftnotesfalse} in \texttt{main.tex}.
    \item Reproduce the no-task-mass cold-start, parcel-prediction, online-inertia and drop-regression batches from frozen logs and add confidence intervals.
    \item Run zero-floor, no-ratchet and envelope-free ablations; state explicitly which mass-valued safety settings remain in each arm.
    \item Verify that no run passes parcel truth, empty-mass initialisation or the former $0.95m_b$ lower bound to MHE.
    \item Verify every numerical claim against a retained log or analysis script.
    \item Run BibTeX and compile LaTeX until all cross-references are resolved.
    \item Check the final PDF against the current NTU MAE formatting checklist.
\end{enumerate}
\end{spacing}

\end{appendices}
